\documentclass[11pt,a4paper]{article}

% Pakete
\usepackage[utf8]{inputenc}
\usepackage[T1]{fontenc}
\usepackage[ngerman]{babel}
\usepackage[left=18mm,right=18mm,top=18mm,bottom=18mm]{geometry}
\usepackage{helvet}
\renewcommand{\familydefault}{\sfdefault}
\usepackage{amsmath,amssymb}
\usepackage{tikz}
\usepackage{tcolorbox}
\tcbuselibrary{skins,breakable}
\usepackage{enumitem}
\usepackage{tabularx}
\usepackage{siunitx}
\sisetup{locale=DE, per-mode=fraction}
\usepackage{colortbl}

% Fachfarbe
\definecolor{physik}{HTML}{2c5aa0}
\definecolor{physiklight}{HTML}{e3f2fd}
\definecolor{lehrer}{HTML}{6a1b9a}
\definecolor{lehrerlight}{HTML}{f3e5f5}

% Box-Stile
\tcbset{
    lehrerbox/.style={
        colback=lehrerlight,
        colframe=lehrer,
        fonttitle=\bfseries,
        title=#1,
        boxrule=1pt,
        arc=2pt,
        left=6pt, right=6pt, top=5pt, bottom=5pt
    },
    infobox/.style={
        colback=physiklight,
        colframe=physik,
        boxrule=1pt,
        arc=2pt,
        left=6pt, right=6pt, top=5pt, bottom=5pt
    }
}

% Kopfzeile
\usepackage{fancyhdr}
\pagestyle{fancy}
\fancyhf{}
\lhead{\small Physik 12 GK}
\chead{\small \textcolor{lehrer}{LEHRERHINWEISE}}
\rhead{\small Fr 16.01.2026}
\renewcommand{\headrulewidth}{0.4pt}

% Keine Einrückung
\setlength{\parindent}{0pt}
\setlength{\parskip}{6pt}

\begin{document}

% ============================================================
% TITEL
% ============================================================

{\Large\textbf{Photoeffekt Vertiefung + De-Broglie-Wellenlänge}}\\[0.3em]
{\normalsize Doppelstunde (90 min) | Fr 16.01.2026}

\vspace{0.3em}
\hrule height 1.5pt
\vspace{0.5em}

% ============================================================
% ÜBERBLICK
% ============================================================

\begin{tcolorbox}[lehrerbox={Stundenüberblick}]
\textbf{Lernziele:}
\begin{itemize}[nosep]
\item Photogleichung anwenden (Abiturformat)
\item De-Broglie-Hypothese verstehen und anwenden
\item Größenordnungen einschätzen (messbar vs. nicht messbar)
\end{itemize}

\vspace{0.3em}
\textbf{Materialien:}
\begin{itemize}[nosep]
\item AB-03-de-broglie.pdf (1 pro SuS, 2 Seiten)
\item Tablets/Smartphones für SIM-03 und LP-03
\item Beamer für Simulation-Demo
\item TB-03-de-broglie.pdf (optional)
\end{itemize}

\vspace{0.3em}
\textbf{Differenzierung:}
\begin{itemize}[nosep]
\item Pflichtaufgaben 1--3 (24 BE) für alle
\item Zusatzaufgaben Z1--Z3 für Schnelle
\end{itemize}
\end{tcolorbox}

\vspace{0.5em}

% ============================================================
% VERLAUFSPLAN
% ============================================================

\textbf{Stundenverlauf}

\vspace{0.3em}

\begin{tabularx}{\textwidth}{|c|l|X|l|}
\hline
\rowcolor{physiklight}
\textbf{Zeit} & \textbf{Phase} & \textbf{Inhalt} & \textbf{Material} \\
\hline
0--5 & Einstieg & Wiederholung Photogleichung, Überblick DS & Beamer \\
\hline
5--35 & \textbf{Block 1} & \textbf{Photoeffekt Vertiefung} & \\
& & Aufgabe 1: Abituraufgabe Cäsium (selbstständig) & LP + AB \\
& & Kasten 1: Merksatz ergänzen & \\
\hline
35--40 & Pause & 5 min Pause & \\
\hline
40--45 & Input & De-Broglie-Hypothese einführen & Beamer/TB \\
\hline
45--55 & Demo & SIM-03 gemeinsam erkunden & Beamer + SIM \\
\hline
55--80 & \textbf{Block 2} & \textbf{De-Broglie-Rechnungen} & \\
& & Aufgabe 2: Elektron bei 150 V (gemeinsam beginnen) & LP + AB \\
& & Aufgabe 3: Tennisball (selbstständig) & \\
& & Kasten 3: Merksatz ergänzen & \\
\hline
80--90 & Sicherung & Zusammenfassung, Zusatzaufgaben für Schnelle & \\
\hline
\end{tabularx}

\vspace{0.5em}

% ============================================================
% DIDAKTISCHE HINWEISE
% ============================================================

\begin{tcolorbox}[lehrerbox={Didaktische Hinweise}]
\textbf{Block 1 -- Photoeffekt (Cäsium):}
\begin{itemize}[nosep]
\item Material: Cäsium mit $W_A = \SI{1,9}{eV}$ (niedrige Austrittsarbeit, gut im sichtbaren Bereich)
\item Aufgabe 1 ist bewusst im Abiturformat gehalten
\item Wichtig bei 1c: y-Achsenabschnitt = $-W_A$ (negative Energie = Energieschuld)
\item Bei 1d: Ein-Photon-ein-Elektron-Prinzip betonen
\end{itemize}

\vspace{0.3em}
\textbf{Block 2 -- De-Broglie:}
\begin{itemize}[nosep]
\item Simulation als Demo nutzen: verschiedene Objekte durchklicken
\item Bei Aufgabe 2 gemeinsam den ersten Schritt (Geschwindigkeit) rechnen
\item Tennisball-Aufgabe zeigt Grenzen der Messbarkeit: $\lambda \approx 10^{-34}$ m ist 24 Größenordnungen kleiner als Atomabstände
\end{itemize}

\vspace{0.3em}
\textbf{Typische Schülerfehler:}
\begin{itemize}[nosep]
\item Verwechslung von $h$ in J·s und eV·s (beides auf AB angegeben)
\item Einheitenfehler bei Geschwindigkeitsformel
\item Vergessen, dass pm = $10^{-12}$ m
\item Masse des Tennisballs nicht in kg umrechnen (58 g → 0,058 kg)
\end{itemize}
\end{tcolorbox}

\vspace{0.5em}

% ============================================================
% LÖSUNGEN KURZFORM
% ============================================================

\textbf{Lösungen (Kurzform)}

\vspace{0.3em}

\begin{tabularx}{\textwidth}{|l|X|}
\hline
\rowcolor{physiklight}
\textbf{Aufgabe} & \textbf{Lösung} \\
\hline
1a (2 BE) & $f_G = \SI{4,6e14}{Hz}$ (Ablesen aus Diagramm) \\
\hline
1b (3 BE) & $E_{\text{kin}} = \SI{4,14e-15}{eV.s} \cdot \SI{7e14}{Hz} - \SI{1,9}{eV} = \SI{2,9}{eV} - \SI{1,9}{eV} = \SI{1,0}{eV}$ \\
\hline
1c (3 BE) & y-Achsenabschnitt bei $f=0$: $E_{\text{kin}} = -W_A = -\SI{1,9}{eV}$. Physikalisch: Austrittsarbeit \\
\hline
1d (3 BE) & $f = \SI{3e14}{Hz} < f_G$: Photonenenergie $E = \SI{1,24}{eV} < W_A$. Mehr Photonen ändern nichts. \\
\hline
2a (2 BE) & $v = \sqrt{\frac{2 \cdot \SI{1,6e-19}{C} \cdot \SI{150}{V}}{\SI{9,1e-31}{kg}}} = \SI{7,3e6}{m/s}$ \\
\hline
2b (2 BE) & $\lambda = \frac{\SI{6,63e-34}{J.s}}{\SI{9,1e-31}{kg} \cdot \SI{7,3e6}{m/s}} = \SI{100}{pm}$ \\
\hline
2c (2 BE) & $\lambda = 100$ pm $\approx$ Atomabstand ($\approx 200$ pm) $\rightarrow$ Beugung an Kristallen möglich \\
\hline
3a (3 BE) & $p = \SI{0,058}{kg} \cdot \SI{50}{m/s} = \SI{2,9}{kg.m/s}$, $\lambda = \SI{2,3e-34}{m}$ \\
\hline
3b (4 BE) & $\lambda$ ist $10^{24}$-mal kleiner als Atomabstände $\rightarrow$ kein Hindernis für Beugung existiert \\
\hline
\end{tabularx}

\vspace{0.5em}

\textbf{Zusatzaufgaben:}

\begin{tabularx}{\textwidth}{|l|X|}
\hline
Z1 (2 BE) & $\lambda = \SI{7,9e-14}{m} = \SI{0,079}{pm}$ \\
\hline
Z2 (3 BE) & $U = \SI{151}{V}$ (Umstellen: $v = h/(m_e \cdot \lambda)$, dann $U = m_e v^2 / (2e)$) \\
\hline
Z3 (3 BE) & Elektron hat größere $\lambda$, da $p \propto \sqrt{m}$ und $\lambda = h/p$ \\
\hline
\end{tabularx}

\vspace{0.5em}

% ============================================================
% KASTEN-LÖSUNGEN
% ============================================================

\textbf{Lückentexte (Kästen)}

\vspace{0.3em}

\begin{tcolorbox}[infobox]
\textbf{Kasten 1:} Ein \textbf{Photon} überträgt seine Energie auf \textbf{ein} Elektron. Unterhalb der \textbf{Grenzfrequenz} findet kein Photoeffekt statt -- unabhängig von der \textbf{Intensität}.

\vspace{0.3em}
\textbf{Kasten 3:} $\lambda = $ \textbf{h} $/$ \textbf{p}. Je größer die \textbf{Masse / der Impuls}, desto kleiner $\lambda$. Bei makroskopischen Objekten: \textbf{keine Beugung beobachtbar}.
\end{tcolorbox}

\vspace{0.5em}

% ============================================================
% AUSBLICK
% ============================================================

\begin{tcolorbox}[lehrerbox={Ausblick nächste Stunde}]
\textbf{Di 21.01.2026 (DS):} Doppelspaltexperiment mit Elektronen

\begin{itemize}[nosep]
\item Davisson-Germer / Jönsson-Experiment
\item Nachweis der Welleneigenschaften von Elektronen
\item Welle-Teilchen-Dualismus
\end{itemize}

\textbf{Vorbereitung:} Simulation Doppelspalt vorbereiten
\end{tcolorbox}

\end{document}
